\documentclass[12pt,a4paper]{article}
\usepackage[T2A]{fontenc}
\usepackage[utf8]{inputenc}
\usepackage[russian]{babel}
\usepackage{amsmath,amssymb}
\usepackage[margin=2cm]{geometry}
\usepackage{listings}
\usepackage{xcolor}
\usepackage{hyperref}

\definecolor{titleblue}{RGB}{26,35,126}
\definecolor{codebg}{RGB}{30,30,30}
\definecolor{codefg}{RGB}{212,212,212}
\definecolor{kw}{RGB}{86,156,214}
\definecolor{fn}{RGB}{220,220,170}
\definecolor{cm}{RGB}{106,153,85}
\definecolor{st}{RGB}{206,145,120}

\lstdefinestyle{py}{
  language=Python,
  basicstyle=\ttfamily\small\color{codefg},
  backgroundcolor=\color{codebg},
  keywordstyle=\color{kw}\bfseries,
  commentstyle=\color{cm}\itshape,
  stringstyle=\color{st},
  functionnamestyle=\color{fn},
  numbers=none,
  frame=single,
  rulecolor=\color{codebg},
  breaklines=true,
  showstringspaces=false,
  tabsize=4,
  xleftmargin=1cm,
  xrightmargin=1cm,
  aboveskip=6pt,
  belowskip=6pt,
}

\lstset{style=py}

\title{\color{titleblue}\textbf{Описание функций сглаживания\\(Гибридная версия: математика + Python)}}
\author{}
\date{}

\begin{document}
\maketitle

%% === 1. BOXCAR ===
\section{Равномерное скользящее среднее (Boxcar)}

\subsection{Математическая формула}

Равномерное скользящее среднее по окну $M$ (половина окна $m = M // 2$):
\[
\text{start}(i) = \max(0,\; i + 1 - m), \qquad
\text{end}(i) = \min(m + i, \; N),
\]
\[
W(i) = \text{end}(i) - \text{start}(i), \qquad
y[i] = \frac{1}{W(i)} \sum_{j = \text{start}(i)}^{\text{end}(i) - 1} x[j].
\]

На границах массива окно обрезается (clamping).

\subsection{Реализация --- Python}

\begin{lstlisting}
def boxcar(win, data):
    half = win // 2
    n = len(data)
    avg = []
    for i in range(n):
        start = max(0, i + 1 - half)
        end = min(half + i, n)
        sum_array = sum(data[start:end])
        aver = sum_array / len(data[start:end])
        avg.append(aver)
    return avg
\end{lstlisting}

\subsection{Свойства}

\begin{itemize}
  \item Веса: равные $1/W(i)$ для всех элементов в окне.
  \item АЧХ: $H(f) = \sin(\pi f M) \big/ \bigl(M \cdot \sin(\pi f)\bigr)$.
  \item Сложность: $O(N \cdot \text{win})$.
  \item Граничные условия: clamping (обрезка окна).
\end{itemize}

%% === 2. GAUSSIAN ===
\section{Гауссов фильтр (Gaussian)}

\subsection{Математическая формула}

Ядро Гаусса с $\sigma = M / 6$:
\[
w[k] = \exp\!\left(-\frac{(k - m)^2}{2\sigma^2}\right), \quad k = 0, \ldots, M{-}1,
\qquad
w\_norm[k] = \frac{w[k]}{\sum_{j=0}^{M-1} w[j]}.
\]

Сигнал дополняется симметрическим отражением (mirror reflection) на $m$ элементов с каждой стороны, затем выполняется свёртка:
\[
y[i] = \sum_{k=0}^{M-1} w\_norm[k] \cdot x\_padded[i + k].
\]

\subsection{Реализация --- Python}

\begin{lstlisting}
def gaussian(data, win):
    half = win // 2
    sigma = win / 6.0
    n = len(data)
    k = [math.exp(-0.5*(i - half)**2 / sigma**2)
         for i in range(win)]
    k = [x/sum(k) for x in k]
    p = ([data[i] for i in range(half-1,-1,-1)]
         + data
         + [data[i] for i in range(n-1,n-half-1,-1)])
    return [sum(k[j]*p[i+j] for j in range(win))
            for i in range(n)]
\end{lstlisting}

\subsection{Свойства}

\begin{itemize}
  \item Веса: убывают по закону Гаусса, нет артефактов Гиббса.
  \item Симметричность ядра: нулевая фаза (сдвиг отсутствует).
  \item $\sigma = M/6$: $\sim$99.7\% энергии внутри окна.
  \item Граничные условия: mirror reflection (отражение).
  \item Сложность: $O(N \cdot \text{win})$.
\end{itemize}

%% === 3. PROGRESSIVE ===
\section{Прогрессивное среднее (Progressive)}

\subsection{Математическая формула}

\[
y[i] = \frac{1}{i + 1} \sum_{k=0}^{i} x[k], \qquad i = 0, 1, \ldots, N{-}1.
\]

Рекуррентно: $S[i] = S[i{-}1] + x[i]$, \quad $y[i] = S[i] \big/ (i + 1)$.

\subsection{Реализация --- Python}

\begin{lstlisting}
def progressive(data):
    result = []
    total = 0.0
    for i, value in enumerate(data):
        total += value
        result.append(total / (i + 1))
    return result
\end{lstlisting}

\subsection{Свойства}

\begin{itemize}
  \item Размер окна: $i + 1$ (растёт от $1$ до $N$).
  \item Дисперсия: $O(1/i)$ --- минимум на последнем отсчёте.
  \item Сложность: $O(N)$.
\end{itemize}

%% === 4. HYBRID ===
\section{Гибридное сглаживание (Hybrid)}

\subsection{Математическая формула}

\[
y[i] =
\begin{cases}
y_{\text{boxcar}}[i], & 0 \leq i < M, \\[4pt]
(1 - \lambda) \cdot y_{\text{boxcar}}[i] + \lambda \cdot y_{\text{prog}}[i], & M \leq i < 2M, \\[4pt]
y_{\text{prog}}[i], & i \geq 2M,
\end{cases}
\]
где $\lambda(i) = (i - M) / M$, \quad $\lambda \in [0, 1)$.

\subsection{Реализация --- Python}

\begin{lstlisting}
def hybrid(data, win):
    n = len(data)
    box = boxcar(win, data)
    prog = progressive(data)

    out = []
    for i in range(n):
        if i < win:
            out.append(box[i])
        elif i < 2 * win:
            z = (i - win) / win
            k = (1.0 - z) * box[i] + z * prog[i]
            out.append(k)
        else:
            out.append(prog[i])
    return out
\end{lstlisting}

\subsection{Зоны}

\begin{center}
\begin{tabular}{|l|l|l|l|}
\hline
\textbf{Зона} & \textbf{Диапазон} & \textbf{Формула} & \textbf{Поведение} \\
\hline
I   & $0 \leq i < M$     & $y = y_{\text{boxcar}}$   & Стабильное усреднение \\
II  & $M \leq i < 2M$    & $(1{-}\lambda) \cdot box + \lambda \cdot prog$ & Линейный crossfade \\
III & $i \geq 2M$        & $y = y_{\text{prog}}$      & Минимальная дисперсия \\
\hline
\end{tabular}
\end{center}

\subsection{Свойства}

\begin{itemize}
  \item Коэффициент интерполяции: $\lambda(i) = (i - M) / M$, $\lambda \in [0, 1)$.
  \item Непрерывность: гарантирована на всех границах зон.
  \item Переходная зона: ширина $M$ отсчётов.
  \item Сглаживание: монотонно усиливается от зоны I к зоне III.
\end{itemize}

\end{document}
